\documentclass[11pt,a4paper]{article}

% ── Packages ──────────────────────────────────────────────────────────
\usepackage[margin=2.5cm]{geometry}
\usepackage{amsmath,amssymb,amsthm}
\usepackage{algorithm,algpseudocode}
\usepackage{booktabs}
\usepackage{hyperref}
\usepackage{cleveref}
\usepackage{enumitem}
\usepackage{xcolor}
\usepackage{tcolorbox}
\usepackage{graphicx}

\tcbuselibrary{theorems,skins,breakable}

\newtcolorbox{yourcode}{
  colback=gray!5, colframe=gray!60, fonttitle=\bfseries,
  title=Your current code (core\_ops.jl\,:\,251), breakable
}
\newtcolorbox{idea}[1]{
  colback=blue!3, colframe=blue!50!black, fonttitle=\bfseries,
  title={#1}, breakable
}
\newtcolorbox{verdict}{
  colback=green!5, colframe=green!40!black, fonttitle=\bfseries,
  title={Final ranking for BRIDGE transfer mode}, breakable
}

\newcommand{\Trelax}{T_{\mathrm{relax}}}
\newcommand{\ReLU}{\mathrm{ReLU}}
\newcommand{\score}{\mathrm{score}}
\newcommand{\gap}{\mathrm{gap}}
\newcommand{\widthop}{\mathrm{width}}

\title{Smarter Relaxation--Threshold Selection\\for the BRIDGE Conditional-Triangle Relaxation\\[6pt]
\large A survey of methods from the NN-verification literature\\and their adaptation to your transfer-mode MIP}
\author{Research notes for Dana / VHAGaR}
\date{March 2026}

% ======================================================================
\begin{document}
\maketitle
\tableofcontents
\newpage

% ======================================================================
\section{Your current setting}\label{sec:setting}
% ======================================================================

\subsection{The verification problem}

You have two neural networks $N_1$ (the reference) and $N_2$ (the candidate).
In \textbf{transfer mode} you build a single MIP that encodes four forward
passes---$N_1(x)$, $N_2(x)$, $N_1(x')$, $N_2(x')$---and solve for the
minimum perturbation $\delta$ such that $N_2$ misclassifies.

Every \emph{unstable} ReLU neuron (one whose pre-activation interval straddles
zero, $l<0<u$) normally requires a binary variable~$a$ and the standard
big-$M$ encoding:
\begin{align}
  \hat z &\ge 0, \qquad \hat z \ge z, \label{eq:relu-lb}\\
  \hat z &\le z - l\,(1 - a), \qquad \hat z \le u\,a. \label{eq:relu-ub}
\end{align}
Binary variables are the main source of MIP complexity.

\subsection{The conditional-triangle relaxation (BRIDGE eqs.~4 and 6)}

For neurons in the $N_2$ passes (\texttt{n2\_org} and \texttt{n2\_pert}), you
can \emph{eliminate} the binary variable~$a_{N_2}$ by \textbf{conditioning on}
$N_1$'s binary $a_{N_1}$ for the same neuron.  Two sub-cases arise:

\begin{itemize}[leftmargin=2em]
  \item \textbf{Active case} ($a_{N_1}=1$): the pre-activation lives in the
        ``active conditional interval'' $[l_A,\; u_A]$ where
        \[
          l_A = l_{\mathrm{int}}, \qquad
          u_A = u_{\mathrm{pre}} + u_{\mathrm{int}}.
        \]
  \item \textbf{Inactive case} ($a_{N_1}=0$): the pre-activation lives in the
        ``inactive conditional interval'' $[l_I,\; u_I]$ where
        \[
          l_I = l_{\mathrm{pre}} + l_{\mathrm{int}}, \qquad
          u_I = u_{\mathrm{int}}.
        \]
\end{itemize}

Here $[l_{\mathrm{int}}, u_{\mathrm{int}}]$ are the \emph{diff bounds}
($z_{N_2} - z_{N_1}$, for the \texttt{n2\_org} pass) or \emph{composed bounds}
($z_{N_2\text{pert}} - z_{N_1}$, for \texttt{n2\_pert}), and
$[l_{\mathrm{pre}}, u_{\mathrm{pre}}]$ are $N_1$'s pre-activation bounds.

A standard triangle relaxation is applied inside each conditional case using
big-$M$ switching on $a_{N_1}$.  The binary $a_{N_2}$ is \textbf{not created}.

\subsection{The current threshold decision}

\begin{yourcode}
\begin{verbatim}
int_width = u_int - l_int
if int_width < relaxation_threshold      # default: 0.5
    # → relax (eliminate binary)
else
    # → keep exact (standard big-M with binary)
end
\end{verbatim}
\end{yourcode}

\noindent\textbf{Problem.}
This is a single, hand-tuned scalar applied identically to every neuron in every
layer.  An interval width of 0.3 may be ``tight'' in one layer (where
$u - l = 50$) but ``loose'' in another (where $u - l = 0.4$).  There is no
theoretical justification for any particular default value.


% ======================================================================
\section{Method 1: Normalised interval width}\label{sec:method1}
% ======================================================================

\begin{idea}{Method 1 --- Normalised interval width (trivial change)}
Replace the absolute width with a \emph{relative} width, normalised by the
neuron's own pre-activation range.
\end{idea}

\subsection{Definition}
Instead of
$\widthop = u_{\mathrm{int}} - l_{\mathrm{int}}$,
compute
\begin{equation}\label{eq:relwidth}
  \widthop_{\mathrm{rel}}
  \;=\;
  \frac{u_{\mathrm{int}} - l_{\mathrm{int}}}{u - l},
\end{equation}
where $[l,u]$ is the neuron's own pre-activation interval (the one used for the
standard big-$M$ encoding).  The threshold $\Trelax$ now lives in $[0,1]$ and is
unitless.

\subsection{Intuition}
Consider two neurons:
\begin{center}
\begin{tabular}{lcccc}
\toprule
Neuron & $u-l$ & $u_{\mathrm{int}}-l_{\mathrm{int}}$ & $\widthop_{\mathrm{rel}}$ & safe to relax? \\
\midrule
A & 100 & 0.4 & 0.004 & \textcolor{green!60!black}{yes --- interval is tiny relative to range} \\
B & 0.5 & 0.4 & 0.8   & \textcolor{red!70!black}{no --- almost the full range is uncertain} \\
\bottomrule
\end{tabular}
\end{center}
The absolute threshold $\Trelax=0.5$ relaxes both.  The relative threshold
correctly distinguishes them.

\subsection{Implementation}
One-line change in \texttt{core\_ops.jl:251}:
\begin{verbatim}
  # OLD:
  if int_width < relaxation_threshold

  # NEW:
  if int_width / (u - l) < relaxation_threshold
\end{verbatim}

\subsection{Limitations}
This is better but still a single scalar threshold.  It does not account for
\emph{how much formulation tightness} you lose when you relax a specific neuron,
nor does it consider the neuron's influence on the objective.


% ======================================================================
\section{Method 2: Triangle relaxation--gap scoring}\label{sec:method2}
% ======================================================================

\begin{idea}{Method 2 --- Relaxation gap area}
Score each neuron by the \emph{area} of the over-approximation triangle
introduced by the relaxation.
\end{idea}

\subsection{Background: the triangle relaxation gap}

When you relax a split ReLU on $[l,u]$ with $l<0<u$, the convex hull of
$\hat z = \ReLU(z)$ is exactly:
\[
  \hat z \ge 0, \quad \hat z \ge z, \quad
  \hat z \le \frac{u}{u-l}(z - l).
\]
The \emph{gap} between the ReLU graph and the upper facet of the triangle is
a right triangle with area
\begin{equation}\label{eq:gap}
  \gap(l,u) \;=\; \frac{u \cdot |l|}{2\,(u - l)}.
\end{equation}
This is maximised when $u \approx |l|$ (symmetric split) and goes to zero when
either $l \to 0^{-}$ or $u \to 0^{+}$ (near-stable neuron).

\subsection{Adaptation to conditional triangles}

In your setting, the relaxation is \emph{conditional} on $a_{N_1}$.
Each case has its own triangle.  The total relaxation looseness is bounded by
the worse of the two:
\begin{equation}\label{eq:condgap}
  \gap_{\mathrm{cond}} \;=\;
  \max\!\Big(\,
    \gap_A(l_A, u_A),\;
    \gap_I(l_I, u_I)
  \,\Big),
\end{equation}
where $\gap_A$ and $\gap_I$ use \cref{eq:gap} applied to the active and
inactive conditional intervals respectively, with the convention that
$\gap(l,u) = 0$ if $l \ge 0$ (entirely active) or $u \le 0$ (entirely
inactive).

\subsection{Decision rule}
\begin{equation}
  \text{relax neuron $(m,k)$}
  \quad\Longleftrightarrow\quad
  \gap_{\mathrm{cond}}(m,k) < \Trelax.
\end{equation}

\subsection{Why this is better than width}
\begin{itemize}
  \item It directly measures the \emph{volume of over-approximation} you
        introduce --- the thing that actually weakens your MIP bound.
  \item A neuron with a wide interval but where one conditional case is entirely
        active ($l_A \ge 0$) gets a smaller gap score, because that case
        contributes zero gap.
  \item It naturally handles asymmetric intervals ($u \gg |l|$ or vice versa),
        where the width is large but the gap is small.
\end{itemize}

\subsection{Implementation}
\begin{verbatim}
  function tri_gap(l, u)
      if l >= 0.0 || u <= 0.0
          return 0.0
      end
      return u * (-l) / (2.0 * (u - l))
  end

  gap_A = tri_gap(lA, uA)
  gap_I = tri_gap(lI, uI)
  cond_gap = max(gap_A, gap_I)

  if cond_gap < relaxation_threshold
      # relax
  end
\end{verbatim}

\subsection{Limitations}
Still uses a fixed threshold (though now with a geometrically meaningful unit).
Does not account for how important this neuron is to the \emph{objective}.


% ======================================================================
\section{Method 3: Budget-based with layer-adaptive percentiles}%
\label{sec:method3}
% ======================================================================

\begin{idea}{Method 6 --- Layer-adaptive percentile threshold}
Compute a per-layer threshold as a percentile of that layer's interval widths,
so each layer relaxes the same \emph{fraction} of its neurons.
\end{idea}

\subsection{Definition}

During \texttt{compute\_diff\_and\_comp\_bounds()}, after computing interval
bounds for each ReLU layer~$m$, collect all widths:
\[
  W_m = \bigl\{\, u_{\mathrm{int}}^{(m,k)} - l_{\mathrm{int}}^{(m,k)}
         \;\big|\; k = 1,\dots,n_m \,\bigr\}.
\]
Set the layer-specific threshold as a quantile:
\begin{equation}
  \Trelax^{(m)} = \mathrm{quantile}(W_m,\; p), \qquad p \in [0,1].
\end{equation}
For example, $p = 0.3$ means: in each layer, relax the 30\% of neurons with
the narrowest intervals.

\subsection{Advantages}
\begin{itemize}
  \item Handles the layer-to-layer scale variation that defeats a single
        global threshold.
  \item Gives you a direct knob: the \emph{fraction} $p$ of neurons to relax
        per layer, which maps to a predictable reduction in binary count.
  \item Simple to implement inside your existing
        \texttt{compute\_diff\_and\_comp\_bounds()} function.
\end{itemize}

\subsection{Limitations}
\begin{itemize}
  \item Still uses a single scalar $p$ (now a percentile rather than a width).
  \item Not objective-aware: does not consider how important each neuron is
        to the bound.
  \item Treats every layer equally (same fraction $p$), which may not be
        optimal---deeper layers often have wider intervals but lower influence.
\end{itemize}


% ======================================================================
\section{Method 4: Deriving $\Trelax$ from MIPGap (no hyperparameter)}%
\label{sec:method4}
% ======================================================================

\begin{idea}{Method 7 --- MIPGap-derived threshold: compute $\Trelax$ from
quantities you already have}
Instead of choosing $\Trelax$ as a hyperparameter, \textbf{derive it} from
(i)~the MIP optimality tolerance you already set, and (ii)~the PGD attack
objective you already compute.  The result: a per-instance, fully automatic
threshold with no tuning.
\end{idea}

\subsection{The core question}

When you relax a neuron (replace its binary with a conditional triangle), you
make the MIP formulation \emph{looser}.  The LP relaxation bound moves further
from the true optimum.  The question is:

\begin{quote}
\emph{How much looseness can I tolerate before it affects my answer?}
\end{quote}

The answer is already in your Gurobi configuration.

\subsection{Step 1: What MIPGap means}\label{sec:m7:mipgap}

You set \texttt{MIPGap = 0.01} in your MIP solver (see \texttt{mip.jl}).  This
tells Gurobi: ``stop when the relative gap between the best known feasible
solution (the \emph{incumbent}) and the best possible bound is less than 1\%.''

Formally, Gurobi stops when:
\begin{equation}\label{eq:mipgap}
  \frac{|\text{BestBound} - \text{Incumbent}|}{|\text{Incumbent}|} \;\le\; 0.01.
\end{equation}

This means \textbf{your answer already tolerates a 1\% error}.  Any relaxation
that introduces less than 1\% degradation in the bound is invisible --- it
cannot change the outcome.

\subsection{Step 2: What the PGD attack gives you}\label{sec:m7:pgd}

Before solving the MIP, your hyper-attack (PGD) finds a feasible adversarial
example and reports \texttt{best\_val}.  For example, from your log:
\[
  \texttt{best\_val} = 46.99.
\]
This is a \emph{lower bound} on the true MIP optimum (for a maximisation
problem; it is the incumbent).  Gurobi uses it as a warm-start.

\subsection{Step 3: The allowed total degradation}\label{sec:m7:total}

Combining Steps 1 and 2, the maximum amount by which the bound can degrade
\emph{without affecting the answer} is:
\begin{equation}\label{eq:allowed}
  \Delta_{\mathrm{allowed}}
  \;=\;
  \texttt{MIPGap} \;\times\; |\texttt{best\_val}|
  \;=\;
  0.01 \times 46.99
  \;=\;
  0.47.
\end{equation}

If the sum of all relaxation-induced looseness is less than $0.47$, the
relaxations \textbf{cannot change the solver's answer} (it would have stopped
at the same gap anyway).

\subsection{Step 4: How one relaxed neuron degrades the bound}\label{sec:m7:one}

When you replace the exact ReLU encoding (with binary variable) by a triangle
relaxation, you introduce an over-approximation.  For a single neuron with
pre-activation bounds $[l, u]$ where $l < 0 < u$, the standard triangle
relaxation has a gap (the area between the true ReLU and the convex upper
envelope):
\begin{equation}\label{eq:gap-reminder}
  g(l, u) \;=\; \frac{u \cdot |l|}{2\,(u - l)}.
\end{equation}

\textbf{Why area?}  The gap area represents the ``volume of over-approximation''
--- the set of $(\hat z, z)$ points that the relaxation considers feasible but
that violate the true ReLU.  A larger area means the LP has more room to
``cheat'' and produce a bound that is further from the truth.

In your conditional-triangle setting, each relaxed $N_2$ neuron has \emph{two}
sub-triangles (active/inactive cases conditioned on $a_{N_1}$).  The
worst-case gap is:
\begin{equation}\label{eq:condgap-m7}
  g_{\mathrm{cond}}(m,k)
  \;=\;
  \max\!\bigl(\,g(l_A, u_A),\;\; g(l_I, u_I)\,\bigr),
\end{equation}
where:
\begin{align}
  l_A &= l_{\mathrm{int}},          & u_A &= u_{\mathrm{pre}} + u_{\mathrm{int}},
  \label{eq:active-bounds}\\
  l_I &= l_{\mathrm{pre}} + l_{\mathrm{int}}, & u_I &= u_{\mathrm{int}},
  \label{eq:inactive-bounds}
\end{align}
which you already compute at \texttt{core\_ops.jl:267--268}.  If a sub-case is
entirely active ($l \ge 0$) or entirely inactive ($u \le 0$), its gap is zero.

\subsection{Step 5: From total budget to per-neuron budget}\label{sec:m7:per}

The total bound degradation from relaxing a set $\mathcal{R}$ of neurons is
bounded by:
\begin{equation}\label{eq:sumgap}
  \text{Total degradation}
  \;\le\;
  \sum_{(m,k) \in \mathcal{R}} g_{\mathrm{cond}}(m,k).
\end{equation}

\textbf{Note:} this is a \emph{conservative upper bound}.  The actual
degradation is typically much smaller because:
\begin{enumerate}
  \item Not all gaps propagate to the objective (many cancel or are attenuated
        by subsequent layers).
  \item Gurobi's cutting planes and presolve recover some of the looseness.
\end{enumerate}

To stay within the budget $\Delta_{\mathrm{allowed}}$, we need:
\[
  \sum_{(m,k) \in \mathcal{R}} g_{\mathrm{cond}}(m,k)
  \;\le\;
  \Delta_{\mathrm{allowed}}.
\]

The simplest (most conservative) way to enforce this is a \textbf{uniform
per-neuron budget}:
\begin{equation}\label{eq:perbudget}
  \boxed{\;
    g_{\mathrm{budget}}
    \;=\;
    \frac{\Delta_{\mathrm{allowed}}}{N_{\mathrm{unstable}}}
    \;=\;
    \frac{\texttt{MIPGap} \times |\texttt{best\_val}|}{N_{\mathrm{unstable}}}
  \;}
\end{equation}
where $N_{\mathrm{unstable}}$ is the total number of unstable $N_2$ neurons
(those with $l < 0 < u$).

\subsection{Step 6: The decision rule}\label{sec:m7:rule}

Replace the current threshold test with:
\begin{equation}\label{eq:m7rule}
  \text{Relax neuron $(m,k)$}
  \quad\Longleftrightarrow\quad
  g_{\mathrm{cond}}(m,k) \;<\; g_{\mathrm{budget}}.
\end{equation}

\textbf{In words:} relax a neuron if and only if the gap it introduces is
smaller than its ``fair share'' of the total allowed degradation.

\subsection{Worked example}\label{sec:m7:example}

From your log file (\texttt{Phase 1 (standard)\_trans(1,1).log}):
\begin{itemize}
  \item \texttt{best\_val} $= 46.99$ (PGD warm-start objective)
  \item \texttt{MIPGap} $= 0.01$ (your Gurobi setting)
  \item Network: CNN1 with $775$ binary variables (unstable ReLU neurons)
  \item In transfer mode, $N_2$ has roughly $N_{\mathrm{unstable}} \approx 775$
        neurons per pass, so $\sim 1550$ total for \texttt{n2\_org} +
        \texttt{n2\_pert}
\end{itemize}

\noindent
The computation:
\begin{align}
  \Delta_{\mathrm{allowed}}
    &= 0.01 \times 46.99 = 0.4699, \\
  g_{\mathrm{budget}}
    &= \frac{0.4699}{1550} \approx 3.03 \times 10^{-4}.
\end{align}

So a neuron is relaxed if its conditional gap is less than $\sim 0.0003$.
This is a \emph{very conservative} threshold (it guarantees the total
degradation is within MIPGap even in the worst case).

\medskip
\noindent
\textbf{If this is too conservative} (too few neurons relaxed), you can use a
\emph{relaxation factor} $\alpha > 1$:
\begin{equation}\label{eq:alpha}
  g_{\mathrm{budget}}
  \;=\;
  \alpha \;\cdot\;
  \frac{\texttt{MIPGap} \times |\texttt{best\_val}|}{N_{\mathrm{unstable}}}.
\end{equation}
Setting $\alpha = 10$ means you accept up to $10\times$ the MIPGap in
relaxation-induced looseness.  Since the sum-of-gaps bound
(\cref{eq:sumgap}) is very conservative (the real degradation is much less),
$\alpha \in [5, 50]$ is reasonable in practice.  The key point is that even
with $\alpha$, the threshold \textbf{adapts to the instance}
(different \texttt{best\_val}, different network size $\to$ different
threshold), unlike the fixed $\Trelax = 0.5$.

\subsection{Implementation}\label{sec:m7:impl}

\textbf{Step A.}  Before building the MIP, compute the budget.
In \texttt{run.jl}, after the PGD hyper-attack returns \texttt{best\_val}:

\begin{verbatim}
  # Derived relaxation budget (no hyperparameter)
  MIPGap      = 0.01                    # your existing setting
  alpha       = 10.0                    # relaxation factor (conservative)
  best_val    = ...                     # from hyper-attack
  N_unstable  = ...                     # count after interval propagation

  gap_budget  = alpha * MIPGap * abs(best_val) / max(1, N_unstable)
\end{verbatim}

\textbf{Step B.}  Count $N_{\mathrm{unstable}}$ during interval propagation.
In \texttt{compute\_diff\_and\_comp\_bounds()}, after each ReLU layer:

\begin{verbatim}
  # Count unstable neurons (l < 0 < u) for N2
  n2_pre_down = n1_pre_down_cur .+ diff_down  # N2 preact lower
  n2_pre_up   = n1_pre_up_cur   .+ diff_up    # N2 preact upper
  N_unstable += count(i -> n2_pre_down[i] < 0 < n2_pre_up[i],
                      eachindex(n2_pre_down))
\end{verbatim}

\textbf{Step C.}  Replace the threshold test in \texttt{core\_ops.jl:251}.
The values \texttt{lA, uA, lI, uI} are already computed at lines 267--268:

\begin{verbatim}
  # OLD (line 251):
  if int_width < relaxation_threshold

  # NEW:
  function tri_gap(l, u)
      (l >= 0.0 || u <= 0.0) && return 0.0
      return u * (-l) / (2.0 * (u - l))
  end

  gap_A = tri_gap(lA, uA)
  gap_I = tri_gap(lI, uI)
  if max(gap_A, gap_I) < gap_budget
\end{verbatim}

\subsection{Why this works}\label{sec:m7:why}

\begin{enumerate}
  \item \textbf{Fully automatic.}  $g_{\mathrm{budget}}$ is computed from
        \texttt{MIPGap} (which you already set), \texttt{best\_val} (which the
        PGD attack already produces), and $N_{\mathrm{unstable}}$ (which you
        can count during interval propagation).  No new hyperparameter.

  \item \textbf{Instance-adaptive.}  Different source/target pairs have
        different \texttt{best\_val} and different numbers of unstable neurons.
        The budget automatically adjusts: harder instances (larger
        \texttt{best\_val}) get a larger budget; bigger networks get a smaller
        per-neuron budget.

  \item \textbf{Theoretically grounded.}  The derivation follows from the
        definition of MIPGap and the geometry of the triangle relaxation.
        The bound in \cref{eq:sumgap} is standard in the MIP literature.

  \item \textbf{Tuning $\alpha$ is robust.}  Unlike $\Trelax$ (which has
        unintuitive units and varies wildly across instances), $\alpha$ is
        a dimensionless multiplier on a conservatism factor.  Any value in
        $[1, 100]$ is reasonable; the results are not sensitive to it because
        the sum-of-gaps bound is already very conservative.
\end{enumerate}

\subsection{Comparison with the current approach}

\begin{center}
\renewcommand{\arraystretch}{1.3}
\begin{tabular}{@{}l c c@{}}
\toprule
& \textbf{Current} ($\Trelax = 0.5$)
& \textbf{Method 7} ($g_{\mathrm{budget}}$) \\
\midrule
What is thresholded?
  & interval width & conditional gap area \\
Units
  & pre-activation units (layer-dependent)
  & objective units (instance-adaptive) \\
Adapts to instance?
  & no (fixed for all instances)
  & yes (via \texttt{best\_val}, $N_{\mathrm{unstable}}$) \\
Adapts across layers?
  & no (same threshold for all layers)
  & yes (gap normalises for scale) \\
Hyperparameter?
  & $\Trelax$ (hard to tune)
  & $\alpha$ (robust, optional) \\
Theoretical basis?
  & none
  & MIPGap tolerance + triangle geometry \\
\bottomrule
\end{tabular}
\end{center}


% ======================================================================
\section{Ranking and recommendation}\label{sec:ranking}
% ======================================================================

\begin{verdict}
Below I rank all four methods for your specific setting: a BRIDGE-style
transfer verifier with Gurobi as the MIP solver, targeting MNIST-scale networks
(CNN1/CNN2 with $\sim$775 binary variables per pass).

\bigskip

\renewcommand{\arraystretch}{1.3}
\begin{tabular}{@{}c l c c c c@{}}
\toprule
\textbf{Rank} & \textbf{Method} & \textbf{Effort} &
  \textbf{Theory} & \textbf{Impact} & \textbf{Eliminates $\Trelax$?} \\
\midrule
1 & MIPGap-derived (\S\ref{sec:method4})
  & \textcolor{green!50!black}{low} & \textcolor{green!50!black}{strong}
  & \textcolor{green!50!black}{high} & yes (fully) \\
2 & Relaxation gap area (\S\ref{sec:method2})
  & low & \textcolor{blue!70}{moderate}
  & \textcolor{blue!70}{medium} & no \\
3 & Layer-adaptive percentile (\S\ref{sec:method3})
  & low & \textcolor{orange!80!black}{weak}
  & \textcolor{blue!70}{medium} & no \\
4 & Normalised width (\S\ref{sec:method1})
  & trivial & \textcolor{orange!80!black}{weak}
  & \textcolor{orange!80!black}{low--medium} & no \\
\bottomrule
\end{tabular}

\bigskip

\textbf{Detailed justification:}

\begin{enumerate}[leftmargin=2em, label=\textbf{\arabic*.}]

\item \textbf{MIPGap-derived threshold (rank 1).}
This is the best approach because:
\begin{itemize}
  \item \textbf{Zero new hyperparameters.}  It derives $\Trelax$ from
        \texttt{MIPGap} and \texttt{best\_val}, both of which you already have.
  \item \textbf{Low effort.}  Requires only adding a \texttt{tri\_gap()}
        function and changing the threshold test in \texttt{core\_ops.jl:251}.
        No extra solver calls, no MIP restructuring.
  \item \textbf{Theoretically grounded.}  The derivation
        (\S\ref{sec:method4}) follows from the geometry of the triangle
        relaxation and the definition of MIPGap.  It is not an ad-hoc
        heuristic.
  \item \textbf{Instance-adaptive.}  Different source/target pairs
        automatically get different thresholds (via different
        \texttt{best\_val} and $N_{\mathrm{unstable}}$).
  \item \textbf{Compatible with your current code structure.}  The
        change is a drop-in replacement for line~251 --- no architectural
        changes needed.
  \item The optional relaxation factor $\alpha$ is robust: any value in
        $[5, 50]$ is reasonable, and results are not sensitive to it.
\end{itemize}

\item \textbf{Relaxation gap area (rank 2).}
A good ``local'' improvement that requires no extra solver calls.  It correctly
captures the geometric looseness introduced by relaxation.  But it is
\emph{not objective-aware}: a neuron with a large gap that happens to have
zero influence on the bound will be needlessly kept exact.
Note that Method~4 already uses the gap area as its scoring quantity --- so
this method is effectively subsumed by Method~4, but is simpler if you want
a quick first step.

\item \textbf{Layer-adaptive percentile (rank 3).}
Fixes the cross-layer calibration issue (which is real) but in a crude way.
It treats all neurons in a layer as equally important, which they are not.
The fraction $p$ is still a hyperparameter, though more interpretable than
the current raw width threshold.

\item \textbf{Normalised width (rank 4).}
Trivial to implement and strictly better than the current absolute width.
But the improvement is modest: normalisation helps but is still just a
one-dimensional heuristic with no optimisation signal.

\end{enumerate}
\end{verdict}

\subsection{Recommended implementation path}

\begin{enumerate}
  \item \textbf{Immediate} (30 minutes): implement \textbf{Method~4}
        (MIPGap-derived threshold).  This is a drop-in replacement for the
        current threshold test.  See \S\ref{sec:m7:impl} for the exact code
        changes.  Start with $\alpha = 10$ and adjust if needed.

  \item \textbf{Validation}: run a few source/target pairs and compare the
        number of relaxed neurons (\texttt{relaxation\_condition\_count}) and
        solve times against the fixed $\Trelax = 0.5$ baseline.

  \item \textbf{If too conservative} (too few neurons relaxed): increase
        $\alpha$.  If too aggressive (bound quality degrades): decrease
        $\alpha$.  Any value in $[5, 50]$ is reasonable.
\end{enumerate}


% ======================================================================
\section*{References}
% ======================================================================

\begin{enumerate}[label={[\arabic*]}, leftmargin=2em]

\item\label{ref:tjeng2019}
  V.~Tjeng, K.~Xiao, and R.~Tedrake.
  ``Evaluating Robustness of Neural Networks with Mixed Integer Programming.''
  \emph{ICLR}, 2019.
  \url{https://github.com/vtjeng/MIPVerify.jl}

\item\label{ref:anderson2020}
  R.~Anderson, J.~Huchette, W.~Ma, C.~Tjandraatmadja, and J.P.~Vielma.
  ``Strong Mixed-Integer Programming Formulations for Trained Neural Networks.''
  \emph{Mathematical Programming}, 2020.
  \url{https://dspace.mit.edu/bitstream/handle/1721.1/131368/10107_2020_1474_ReferencePDF.pdf}

\item\label{ref:bunel2020}
  R.~Bunel, P.~Mudigonda, I.~Turkaslan, P.~Torr, J.~Lu, and A.~Kohli.
  ``Branch and Bound for Piecewise Linear Neural Network Verification.''
  \emph{JMLR}, 21:1--39, 2020.
  \url{https://jmlr.org/papers/volume21/19-468/19-468.pdf}

\item\label{ref:li2023}
  Y.~Li et al.
  ``Neuron importance based verification of neural networks via divide and
  conquer.'' \emph{Neurocomputing}, 2023.
  \url{https://www.sciencedirect.com/science/article/abs/pii/S0925231023011189}

\end{enumerate}

\end{document}